% SHARD-ID: AQM-68-QSA-FINITE-SET-TENSOR-SITES
% SHARD-TITLE: Structureless Finite-Set Tensor-Site Association
% SHARD-SUMMARY: Defines the finite-set tensor-site workflow after adding one finite-dimensional Hilbert carrier or observable algebra per site.
% SHARD-SUMMARY: Forms finite tensor products and tensor-by-identity local observable inclusions by local finite-dimensional derivation.
% SHARD-SUMMARY: Separates this workflow from the formal carrier of AQM-67 and reviews only the anchored uniform-qudit cases.
% SHARD-KEYWORDS: quantum-system association, finite set, tensor product, local observables, qudit anchors, extra data

\section{Structureless Finite-Set Tensor-Site Association}
\label{sec:qsa-finite-set-tensor-sites}

\subsection{Status and Local Anchors}

This is QSA Step 04 from
\path{docs/quantum_system_associations/PLAN.md}.  It is a local
finite-dimensional construction using \texttt{CONVENTIONS.md} conventions
\texttt{(ap)} and \texttt{(ar)}, with AQM-64--AQM-66 as the planning and
catalogue anchors.

The input is not a bare finite set alone.  A bare finite set \(X\) supplies
site labels only.  The tensor-site workflow starts after one adds per-site
quantum data.  The earlier AQM-67 workflow instead returned only the formal
carrier \(\mathbb C\langle X\rangle\) from convention \texttt{(aq)}; it did
not include inner products, adjoints, observable algebras, or tensor factors.

\subsection{Tensor-Site Input Data}

There are two finite-dimensional input versions.

\begin{description}
\item[Hilbert-site input.]
For every \(x\in X\), choose a finite-dimensional complex Hilbert space
\(\mathcal H_x\).  Then set
\[
  \mathcal A_x=\operatorname{End}_{\mathbb C}(\mathcal H_x).
\]

\item[Algebra-site input.]
Alternatively, choose directly a finite-dimensional unital \(*\)-algebra
\(\mathcal A_x\) for every \(x\in X\).  This gives an observable-algebra
workflow, but it does not choose a Hilbert carrier unless a representation is
added separately.
\end{description}

These choices are part of the data.  They are not recovered from the
structureless set \(X\).  An ordering of a finite subset may be used only to
write tensor products; the set itself does not supply a preferred order.

\subsection{Finite Tensor Products}

For \(S\subset X\), define in the Hilbert-site version
\[
  \mathcal H(S)=\bigotimes_{x\in S}\mathcal H_x,\qquad
  \mathcal A(S)=\bigotimes_{x\in S}\mathcal A_x
  \subseteq \operatorname{End}_{\mathbb C}(\mathcal H(S)).
\]
The empty tensor products are
\[
  \mathcal H(\emptyset)=\mathbb C,\qquad
  \mathcal A(\emptyset)=\mathbb C.
\]
In the algebra-site version, keep the same formula for \(\mathcal A(S)\) as
the finite algebraic tensor product over \(\mathbb C\).

If \(X=\{x_1,\ldots,x_n\}\) is listed for display and
\(\dim\mathcal H_{x_i}=d_i\), then
\[
  \dim \mathcal H(X)=\prod_{i=1}^n d_i,
  \qquad
  \dim \mathcal A(X)=\prod_{i=1}^n d_i^2
\]
in the full-matrix Hilbert-site version.  This is the elementary finite tensor
bookkeeping for the chosen factors, not a claim that \(X\) determines the
numbers \(d_i\).

\subsection{Local Observable Inclusions}

If \(S\subset T\subset X\), the tensor-site inclusion is
\[
  \iota_{S,T}:\mathcal A(S)\longrightarrow \mathcal A(T),\qquad
  a\longmapsto a\otimes \mathbf 1_{T\setminus S}.
\]
It preserves multiplication, the unit, and adjoints because each operation is
computed factorwise.  For \(S\subset T\subset U\),
\[
  \iota_{T,U}\circ\iota_{S,T}=\iota_{S,U}
\]
by the associativity of the finite tensor product and by tensoring identities
on exactly the missing factors.  Thus subsets of \(X\), ordered by inclusion,
carry a covariant finite local-observable system after the site algebras have
been chosen.

\subsection{Boundary With AQM-67}

The Step 03 carrier and the Step 04 tensor-site system answer different
questions.  AQM-67 gives
\[
  \mathbb C\langle X\rangle=\bigoplus_{x\in X}\mathbb C e_x
\]
with basis labels but no Hilbert or observable structure.  This shard gives
\[
  (X,\{\mathcal H_x\}_{x\in X})
  \longmapsto
  \bigotimes_{x\in X}\mathcal H_x
\]
or its observable-algebra variant.  The first construction has additive
bookkeeping in the labels; the second has tensor-product bookkeeping for
independent coexisting sites.

Even the one-dimensional-site edge case is not an identification of
workflows.  If all \(\mathcal H_x=\mathbb C\), then
\(\bigotimes_x\mathcal H_x\cong\mathbb C\), while AQM-67 has one formal
generator per element of \(X\).  Any later comparison with
\(\mathbb C\langle X\rangle\) must name the extra choices and the comparison
map.

\subsection{Review Anchors and Caution}

The existing arithmetic shards already contain important tensor-site examples,
but each has extra structure beyond a bare finite set.

AQM-14 starts from a finite point set together with a chosen finite symplectic
target \(V\), a chosen field-label subspace
\(E\leq\operatorname{Map}(X,V)\), and the radical quotient before forming its
Weyl observable algebra.  This is a field-label workflow, not a theorem that
every bare finite set has canonical qudit factors.

AQM-22 proves the uniform prime-field case with \(X=\mathbb F_p\) as a
finite set only after choosing all maps \(E=\operatorname{Map}(X,\mathbb
F_p^2)\) and the Weyl operators on \(\ell^2(\mathbb F_p^X)\).  The same shard
separates this from \(X=\Spec\mathbb F_p\), which is the one-site spectrum.

AQM-24 and AQM-25 use finite-spectrum or quasi-compact-open Zariski locality
after residue fields have supplied local factors
\(\ell^2(\kappa(x))\) and algebras
\(\operatorname{End}(\ell^2(\kappa(x)))\).  Their tensor-by-identity
inclusions are anchors for the local-net formula above, but their sites are
residue-field sites, not arbitrary sites canonically attached to a bare set.

AQM-40 proves the product-field spectrum case
\(\Spec(k^n)\): the spectrum has \(n\) discrete residue-field points and the
resulting residue-qudit arithmetic field is the ambient \(n\)-qudit Weyl
net.  That is a uniform finite-qudit anchor for the particular arithmetic
object \(k^n\), not a universal equivalence between structureless finite sets
and all tensor-site systems.

\subsection{Conclusion}

The QSA Step 04 output is available only for the enhanced input
\[
  (X,\{\mathcal H_x\}_{x\in X})
  \quad\text{or}\quad
  (X,\{\mathcal A_x\}_{x\in X}).
\]
After this data is supplied, finite tensor products and tensor-by-identity
local observable inclusions are formal finite-dimensional bookkeeping.  The
bare set alone still supplies only labels; it does not select the site
Hilbert spaces, observable algebras, Weyl labels, stabilizer data, or
dynamics.
